background modelling:\\
1/ continuum combinatorial background: expoential shape\\
2/ misreconstructed B decays which will normally accumulate on the low mass sideband
as the final states would be missing some signal particles: this is described by an error
function.\\
3/ However we also expect contributions (see background description section) from the
B+ to jpsi(ee)K+ mode which is relatively abundant and will contribute on the right
of the signal peak (as in this case an additional charged track is added to the mass
calculation).
This can be estimated with its shape studied on the same-sign-meson region.
This is described a double crystal ball.\\

Follow up on:\\
1/ Simultaneous fit of the combinatorics exponential component between SS and OS.\\
2/ While extensively testing the fitting model, we realised that in most cases
there was some cross-feed between the exponential shape	and the	error function.
This meant that	the error function would at times take extreme shapes in order
to include the combinatorial component.	Hence a	constraint on the error function
parameters have	been implemented in order to keep this component to represent
the misreconstructed B decays.\\
3/ B+ shape is fit from	the SS events.
